\section{Mesh Neural Cellular Automata}

Niklasson et al. \cite{niklassonMNCA} y Mordvintsev et al. \cite{niklassonMNCA2} \cite{mordvintsevUNCA} introducen el concepto de \textit{Self-Organising Textures}, cuyo objetivo es reproducir la apariencia global de una textura de referencia sin generar copias exactas. Para ello, entrenan un NCA con la misma arquitectura neuronal propuesta en \cite{mordvintsevNCA}, de modo que pueda capturar y generar patrones característicos como tableros de ajedrez, líneas o mosaicos. Aunque este enfoque tiene un carácter demostrativo, constituye una evidencia importante de que los NCA pueden aprender y generalizar comportamientos visuales locales en lugar de memorizar imágenes específicas.

Sobre esta base, Niklasson et al. proponen \textit{Mesh Neural Cellular Automata} (\textbf{MeshNCA}) \cite{meshNCA}, un método diseñado para generar texturas dinámicas en superficies tridimensionales. A diferencia de los experimentos en cuadrículas bidimensionales, MeshNCA extiende el mecanismo de autoorganización a mallas arbitrarias, incluso cuando estas se encuentran en movimiento. El modelo mantiene coherencia espacial en distintos escenarios, muestra robustez frente a perturbaciones externas y permite controlar propiedades como densidad, dirección y rotación de los patrones emergentes. Además, varios MeshNCA pueden injertarse de manera conjunta, produciendo transiciones suaves entre texturas sin discontinuidades.

Sin embargo, este enfoque presenta un reto de escalabilidad. El número de células involucradas suele encontrarse en el rango de $10^4 \sim 10^5$, lo que limita la resolución alcanzable a entre $64 \times 64$ y $256 \times 256$ píxeles. Como el tiempo de entrenamiento escala de forma cuadrática con respecto al número de celdas, el aumento de resolución conlleva un cuello de botella en la propagación de información a lo largo de la malla.

Para superar esta limitación, Niklasson et al. proponen en \cite{meshNCA2} un modelo híbrido que combina NCA con Representaciones Neuronales Implícitas (INR, \textit{Implicit Neural Representations}). La idea central es desacoplar el tamaño de la cuadrícula del NCA de la resolución final, delegando el renderizado detallado a un decodificador ligero denominado \textit{Local Pattern Producing Network} (LPPN). De este modo, el NCA mantiene una cuadrícula explícita y relativamente pequeña que evoluciona de manera local, mientras que la LPPN se encarga de producir el color en cualquier punto del espacio continuo a partir de información interpolada de las celdas vecinas.

En términos formales, para cada punto de muestreo $p$, la LPPN recibe dos entradas: (1) un estado celular promedio $s(p)$, obtenido mediante interpolación de las celdas cercanas en un primitivo geométrico, y (2) un vector de coordenadas locales $u(p)$, que codifica la posición relativa de $p$ dentro de dicho primitivo. Dependiendo de la discretización empleada, estos primitivos pueden ser rectángulos, cubos, triángulos o polígonos generales, definidos por coordenadas $\lambda_i(p)$ que cumplen propiedades como unidad de partición y no negatividad. Para garantizar continuidad espacial, las coordenadas se transforman mediante funciones sinusoidales en cuadrículas cartesianas, o bien en mallas triangulares, obteniendo así campos posicionales suaves.

El resultado es un sistema que conserva las propiedades autoorganizativas de los NCA al tiempo que permite alcanzar mayor resolución sin un aumento proporcional del costo computacional. Desde una perspectiva biológica, el NCA puede interpretarse como un plano morfogenético que define la organización global, mientras que la LPPN añade los detalles finos sin necesidad de simular interacciones locales a gran escala. Este paralelismo, aunque principalmente ilustrativo, refuerza la idea de que los NCA pueden actuar como generadores de estructuras robustas y adaptables, mientras que módulos complementarios aportan precisión y escalabilidad.

A pesar de estas mejoras, una limitación persistente es que cada NCA debe entrenarse de manera específica para una textura o patrón dado. Incluso en escenarios donde múltiples NCA coexisten, cada modelo requiere su propio entrenamiento independiente. Esto ha motivado nuevas propuestas orientadas a la construcción de espacios compartidos de NCA, capaces de generalizar a múltiples texturas dentro de un mismo marco de entrenamiento.